\paragraph*{In short}

\begin{table}[htbp]
\centering\small
\begin{tabularx}{\linewidth}{@{}lXX@{}}
\toprule
\textbf{Concept} & \textbf{Primitive} & \textbf{When to use} \\
\midrule
Predicated execution & \texttt{if} & Data-dependent decisions (bounds, conditions) \\
Structured concurrency & \texttt{inthreads.\allowbreak{}async} & Compile-time thread partitioning (producer/consumer, heterogeneous roles) \\
Inter-region signaling & \texttt{shared event} / \texttt{wait} / \texttt{trigger} & Coordination between concurrent regions \\
Host concurrency & \texttt{parallel.async} / \texttt{stream s} & Multi-kernel overlap, non-blocking launch \\
Persistent scheduling & \texttt{if} + \texttt{foreach} + \texttt{\#} & Fixed grid size, tile striping with padding guard \\
\bottomrule
\end{tabularx}
\end{table}

The 1P1C skeleton above is incomplete: without \texttt{wait} / \texttt{trigger}, the consumer can read before the producer has finished writing. Section~\ref{sec:3-8-5} adds the buffering and event protocols that make the pipeline safe, and walks through the resulting kernel stage by stage.
